\section{曲率}
\label{sec:riem-curvature}

例~\ref{ex:riem-sphere-transport}中，球面上绕一圈平行移动回到起点，向量转过了$\pi/2$；平面上无论怎样绕，向量都不变。这个差别不是坐标造成的，而是两种几何本身不同。本节要找一个逐点定义、与坐标无关的量来检测它。思路是把绕圈缩小到无穷小：沿两个方向先后求协变导数，比较两种次序的结果。得到的曲率张量描述内在几何；最后再把它与曲面在环境空间中“看上去弯曲”的程度区分开。

\subsection{曲率张量与截面曲率}
\label{sec:riem-curvature-tensor}

在欧氏空间中，二阶偏导数可以交换次序。对向量场的协变导数，交换次序的差$\nabla_X\nabla_YZ-\nabla_Y\nabla_XZ$一般不为零，但其中有一部分来自方向场本身：即使在平面上，若$X,Y$不可交换，先沿$X$再沿$Y$与先沿$Y$再沿$X$也不会回到同一点。这部分差别由Lie括号$[X,Y]$描述，应当扣除。

\begin{definition}[黎曼曲率张量]
\label{def:riem-curvature-tensor}
设$\nabla$是$(M,g)$的Levi-Civita联络。对$X,Y,Z\in\mathfrak X(M)$定义
\begin{equation}
R(X,Y)Z=\nabla_X\nabla_YZ-\nabla_Y\nabla_XZ-\nabla_{[X,Y]}Z .
\label{eq:riem-curvature}
\end{equation}
$R$称为$(M,g)$的\term{黎曼曲率张量}（Riemann curvature tensor）。
\end{definition}

文献中曲率的符号约定不统一，有的教材把上式整体取负号。本书采用上式，在这一约定下单位球面的截面曲率为正，后面会核对。

\paragraph{曲率是逐点定义的}
式\eqref{eq:riem-curvature}的三项都含有导数，似乎需要知道$X,Y,Z$在一点附近的取值。实际上$R$对三个变量都按函数线性，所以$(R(X,Y)Z)(p)$只依赖$X_p,Y_p,Z_p$。以$Z$为例，把$fZ$代入，由乘积法则
\[
\nabla_X\nabla_Y(fZ)=(XYf)Z+(Yf)\nabla_XZ+(Xf)\nabla_YZ+f\nabla_X\nabla_YZ .
\]
交换$X,Y$后相减，中间两项抵消，第一项变为$([X,Y]f)Z$；而$\nabla_{[X,Y]}(fZ)=([X,Y]f)Z+f\nabla_{[X,Y]}Z$，恰好消去这一项，剩下$fR(X,Y)Z$。对$X$的函数线性可以类似验证，第三项$\nabla_{[X,Y]}Z$在其中同样负责抵消导数项；对$Y$的情形由$R(X,Y)=-R(Y,X)$得到。因此对点$p$处的切向量$u,v,w$可以写$R(u,v)w\in T_pM$。

\paragraph{坐标表示}
取坐标基场代入，Lie括号项消失，由$\nabla_{\partial_i}\nabla_{\partial_j}\partial_k=\nabla_{\partial_i}\bigl(\sum_m\Gamma^m_{j,k}\partial_m\bigr)$展开，得到$R(\partial_i,\partial_j)\partial_k=\sum_\ell R^\ell_{i,j,k}\partial_\ell$，其中
\begin{equation}
R^\ell_{i,j,k}=\partial_i\Gamma^\ell_{j,k}-\partial_j\Gamma^\ell_{i,k}+\sum_m\bigl(\Gamma^m_{j,k}\Gamma^\ell_{i,m}-\Gamma^m_{i,k}\Gamma^\ell_{j,m}\bigr).
\label{eq:riem-curvature-coordinates}
\end{equation}
曲率由度量的二阶导数决定。若在某组坐标中$\bm G$为常矩阵，则所有Christoffel符号为零，$R=0$。反过来，系数不为零并不说明曲率非零：极坐标中Christoffel符号不为零，平面的曲率仍为零，因为$R$是张量，在直角坐标中为零就在任何坐标中为零。

\paragraph{截面曲率}
$R$有三个输入，信息较多。最常用的标量是在一个二维切平面上读出的曲率。

\begin{definition}[截面曲率]
\label{def:riem-sectional-curvature}
设$u,v\in T_pM$线性无关，它们张成的平面$\Pi$的\term{截面曲率}（sectional curvature）为
\begin{equation}
K(\Pi)=K(u,v)=\frac{g\bigl(R(u,v)v,u\bigr)}{g(u,u)g(v,v)-g(u,v)^2}.
\label{eq:riem-sectional-curvature}
\end{equation}
\end{definition}

分母是$u,v$张成的平行四边形面积的平方，由Cauchy--Schwarz不等式，它在$u,v$线性无关时为正。可以证明分子在把$(u,v)$换成同一平面的另一组基时，与分母乘上相同的因子，所以$K$只依赖平面$\Pi$；这一点用到$g(R(u,v)w,z)$关于$(u,v)$和关于$(w,z)$都反对称的性质，这里不展开。二维流形每点只有一个切平面，截面曲率就是一个函数，称为\term{Gauss曲率}（Gaussian curvature）。一维流形没有二维切平面，式\eqref{eq:riem-curvature}中的$R$恒为零。

\begin{example}[单位球面的曲率]
\label{ex:riem-sphere-curvature}
用例~\ref{ex:riem-christoffel-examples}的系数$\Gamma^\theta_{\phi,\phi}=-\sin\theta\cos\theta$、$\Gamma^\phi_{\theta,\phi}=\Gamma^\phi_{\phi,\theta}=\cot\theta$计算$R(\partial_\theta,\partial_\phi)\partial_\phi$。由式\eqref{eq:riem-curvature-coordinates}，其$\theta$分量为
\begin{equation*}
\begin{aligned}
R^\theta_{\theta,\phi,\phi}&=\partial_\theta\Gamma^\theta_{\phi,\phi}-\partial_\phi\Gamma^\theta_{\theta,\phi}+\sum_m\bigl(\Gamma^m_{\phi,\phi}\Gamma^\theta_{\theta,m}-\Gamma^m_{\theta,\phi}\Gamma^\theta_{\phi,m}\bigr)\\
&=(\sin^2\theta-\cos^2\theta)-0+0-\cot\theta\cdot(-\sin\theta\cos\theta)=\sin^2\theta,
\end{aligned}
\end{equation*}
其中只有$m=\phi$的一项$\Gamma^\phi_{\theta,\phi}\Gamma^\theta_{\phi,\phi}$非零；同样可算出$\phi$分量为零。于是$R(\partial_\theta,\partial_\phi)\partial_\phi=\sin^2\theta\,\partial_\theta$，分子$g(R(\partial_\theta,\partial_\phi)\partial_\phi,\partial_\theta)=\sin^2\theta$，分母$1\cdot\sin^2\theta-0=\sin^2\theta$，所以$K=1$。半径为$a$的球面按同样计算得$K=1/a^2$：球越大，曲率越小，越接近平面。
\end{example}

\paragraph{小回路上的平行移动}
曲率与绕圈搬运的联系可以这样表述。在$p$处取线性无关的$u,v$，在法坐标中沿边长为$s$的小“平行四边形”回路（依次沿$su$、$sv$、$-su$、$-sv$方向）平行移动，记所得的线性映射为$P_s:T_pM\to T_pM$。则当$s\to0$时
\[
P_s(w)=w-s^2R(u,v)w+o(s^2),
\]
首项的符号随绕行方向和曲率符号约定而变。这里只说明机制，不给出证明：每条边上的平行移动可以展开到二阶，四段复合时一阶项相互抵消，二阶项中剩下的正是两种求导次序之差。在单位球面北极取$u=\bm e_1$、$v=\bm e_2$，由球面上$R(u,v)w=g(v,w)u-g(u,w)v$（可由例~\ref{ex:riem-sphere-curvature}的结果连同对称性推出）得$R(u,v)$把$v$映到$u$、把$u$映到$-v$，是切平面内转过直角的线性映射。所以小回路的效应是一个与所围面积$s^2$成比例的小转动。

\paragraph{局部曲率与整体回路}
曲率是局部量，有限回路的总效应还受整体形状影响，不能简单地把一点的曲率乘上面积。一个清楚的例子是锥面。从平面上剪去圆心角为$\alpha$（$0<\alpha<2\pi$）的扇形，把两条剪边粘合，再去掉顶点，得到的曲面每一小块都可以摊平成平面的一块，所以曲率处处为零。然而沿绕顶点一周的回路平行移动，在摊开的扇形中向量保持不变，经过粘合边时却要按扇形张角旋转一次，回到出发点时向量转过角$\alpha$。没有曲率的区域上，绕不可收缩的回路仍可以有非平凡的平行移动；这是拓扑与局部几何共同造成的。

\subsection{内在曲率与嵌入弯曲}
\label{sec:riem-intrinsic-extrinsic}

日常说一个曲面“弯”，通常指它在三维空间中的形状。黎曼曲率衡量的则是曲面内部的长度规则能否摊平成平面。这两件事不同：一张纸可以卷成圆柱而不拉伸、不起皱，纸上的长度与角度保持不变；却无法不拉伸地包住一个球。

\paragraph{平面与圆柱}
取半径为$a$的圆柱面，局部参数化为
\begin{equation*}
\Phi(s,z)=(a\cos(s/a),a\sin(s/a),z),\qquad s\in(0,2\pi a).
\end{equation*}两列切向量为$(-\sin(s/a),\cos(s/a),0)$与$(0,0,1)$，都是单位向量且互相正交，所以诱导度量$\bm G=\bm I$。在这组坐标中Christoffel符号全为零，曲率张量为零，截面曲率$K=0$。就内在几何而言，圆柱面的这一块与平面的一个矩形没有区别。

圆柱在环境中的弯曲由法向量的变化体现。单位法向量为$n(s,z)=(\cos(s/a),\sin(s/a),0)$，沿$s$方向的导数是$\frac1a(-\sin(s/a),\cos(s/a),0)$，大小$1/a$；沿$z$方向不变。这两个数$1/a$与$0$称为曲面在该点的\term{主曲率}（principal curvature），它们描述法向量沿各方向转动的快慢，属于嵌入的性质。对半径为$a$的球面，法向量$n=x/a$沿任何单位切方向的导数大小都是$1/a$，两个主曲率都是$1/a$。

两种曲面的Gauss曲率与主曲率之间有一个关系：对$\mathbb R^3$中的曲面，Gauss曲率等于两个主曲率之积。圆柱给出$\frac1a\cdot0=0$，球面给出$\frac1{a^2}$，与前面的内在计算一致。这一关系称为Gauss绝妙定理，这里不给出证明；它的意义是，主曲率依赖嵌入，它们的乘积却只由诱导度量决定。

\paragraph{为什么要区分}
一维的例子更极端。半径为$a$的圆在平面中明显弯曲，法向偏离切线的程度由曲率$1/a$描述；可它的黎曼曲率恒为零，任何一段弧都与一段直线等距。所以，“法向量在变”“切平面在转”都不能作为内在曲率非零的证据，坐标中度量系数随位置变化也不能。反过来，内在平坦也不意味着嵌入不弯，圆柱就是例子。

这一区分在由采样点估计几何时很重要。第~\ref{chap:manifold-recovery}章用切平面拟合局部点云时，点到切平面的偏离量由法向量的变化，即嵌入弯曲控制，在圆或圆柱上也会出现；它反映的是嵌入的主曲率，而不是内在截面曲率。估计误差的分析须保持这一区分。
